% D-type (binary dihedral) sector: no-E7 exclusion, sterile/Dirac character,
% dark-matter candidacy and ADM-style abundance scenario (Paper XXI, DRAFT)
%% hopfion-framework/preamble.tex
%% Shared preamble for all repository files.
%% \input this at the top of any standalone document.
%% Repository files are designed to be \input into papers directly.

\usepackage{amsmath,amssymb,amsthm}
\usepackage{hyperref}
\usepackage{enumitem}
\usepackage{booktabs}
\usepackage{xcolor}
\usepackage{geometry}
\geometry{margin=2.5cm}

%── Theorem environments ──────────────────────────────────────────────────────
\newtheorem{theorem}{Theorem}[section]
\newtheorem{lemma}[theorem]{Lemma}
\newtheorem{proposition}[theorem]{Proposition}
\newtheorem{corollary}[theorem]{Corollary}
\theoremstyle{definition}
\newtheorem{definition}[theorem]{Definition}
\newtheorem{result}[theorem]{Result}
\theoremstyle{remark}
\newtheorem{remark}[theorem]{Remark}
\newtheorem{observation}[theorem]{Observation}
\newtheorem{conjecture}{Conjecture}
\newtheorem{construction}[theorem]{Construction}
\newtheorem{condition}[theorem]{Condition}

%── Repository metadata commands ─────────────────────────────────────────────
% Usage: \status{published} or \status{open} or \status{conjectured}
\newcommand{\status}[1]{\par\noindent{\color{blue}\textbf{Status:} #1}}
% Usage: \source{Paper I, Theorem thm:scale}
\newcommand{\source}[1]{\par\noindent{\color{darkgray}\textbf{Source:} #1}}
% Usage: \residual{0.013\%} or \residual{exact}
\newcommand{\residual}[1]{\par\noindent{\color{teal}\textbf{Residual:} #1}}
% Usage: \depends{foundations/pentagon, foundations/suppression}
\newcommand{\depends}[1]{\par\noindent{\color{purple}\textbf{Depends on:} #1}}

%── Framework macros ─────────────────────────────────────────────────────────
\newcommand{\vp}{\varphi}
\newcommand{\bst}{\beta^{*}}
\newcommand{\Seff}{S_{\mathrm{eff}}}
\newcommand{\TCMB}{T_{\mathrm{CMB}}}
\newcommand{\Lcond}{\Lambda_{\mathrm{cond}}}
\newcommand{\LUV}{\Lambda_{\mathrm{UV}}}
\newcommand{\LQCD}{\Lambda_{\mathrm{QCD}}}
\newcommand{\MPl}{M_{\mathrm{Pl}}}
\newcommand{\ERy}{E_{\mathrm{Ry}}}
\newcommand{\aem}{\alpha_{\mathrm{em}}}
\newcommand{\sinW}{\sin^{2}\!\theta_{W}}
\newcommand{\vEW}{v_{\mathrm{EW}}}
\newcommand{\twoI}{2I}
\newcommand{\twoT}{2T}
\newcommand{\twoO}{2O}
\newcommand{\twoD}{2D}
\newcommand{\Ehat}{\hat{E}_8}
\newcommand{\QH}{Q_H}
\newcommand{\Ztwo}{\mathbb{Z}_2}
\newcommand{\Zthree}{\mathbb{Z}_3}
\newcommand{\SU}{\mathrm{SU}}
\newcommand{\SO}{\mathrm{SO}}
\newcommand{\cdir}{c_{\mathrm{dir}}}
\newcommand{\Odm}{\Omega_{\mathrm{DM}}}
\newcommand{\Ob}{\Omega_{b}}
\newcommand{\rCMB}{\rho_{\mathrm{CMB}}}
\newcommand{\Lobs}{\Lambda_{\mathrm{obs}}}
\newcommand{\rinfty}{\rho_{\infty}}
\newcommand{\oBD}{\omega_{\mathrm{BD}}}
\newcommand{\xNMC}{\xi_{\mathrm{NMC}}}
\newcommand{\Efb}{E_{\mathrm{fb}}}
\newcommand{\Kfb}{K_{\mathrm{fb}}}
\newcommand{\dimq}[1]{d_{#1}}
\newcommand{\deff}{d_{\mathrm{eff}}}
\newcommand{\Kth}{K_{\mathrm{th}}}
\newcommand{\Ja}{\mathcal{J}_a}
\newcommand{\Jfour}{\mathcal{J}_4}
\newcommand{\Jiso}{\mathcal{J}_{\mathrm{iso}}}
\newcommand{\rhoinf}{\rho_{\infty}}
\newcommand{\rhoa}{\rho_{\mathrm{atom}}}
\newcommand{\Heff}{H_{\mathrm{eff}}}
\newcommand{\epsr}{\varepsilon(r)}
\newcommand{\Hcond}{H_{\mathrm{cond}}}
\newcommand{\eps}{\varepsilon}
\newcommand{\Rzero}{R_0}
\newcommand{\Cstar}{C^*}
\newcommand{\mH}{m_H}
\newcommand{\aem}{\alpha_{\mathrm{em}}}
\newcommand{\ms}{\mu^*}
\newcommand{\lam}{\lambda}
\newcommand{\fzero}{f_0}
\newcommand{\ftwo}{f_2}
\newcommand{\Itor}[1]{I_{#1}^{\mathrm{tor}}}
\newcommand{\Jfb}{\mathcal{J}_{\mathrm{fb}}}
\newcommand{\bz}{\beta_0}
\newcommand{\sopt}{s_{\mathrm{opt}}}
\newcommand{\leff}{\lambda_{\mathrm{eff}}}
\newcommand{\ord}{\operatorname{ord}}
\newcommand{\G}{\mathcal{G}}
\newcommand{\Ia}{I_{2a}}
\newcommand{\Ib}{I_{4}}
\newcommand{\kk}{\kappa}
\newcommand{\aInv}{\alpha^{-1}}
\newcommand{\thetaG}{\theta_{\mathrm{g}}}
\newcommand{\LSC}{\mathrm{LSC}}
\newcommand{\mWZW}{m_{\mathrm{WZW}}}
\newcommand{\tang}{\dot\Gamma}
\newcommand{\Jones}{V}
\newcommand{\lk}{\mathrm{lk}}
\newcommand{\phistar}{\phi^*}
\newcommand{\CC}{\mathbb{C}}
\newcommand{\HH}{\mathbb{H}}
\newcommand{\OO}{\mathbb{O}}
\newcommand{\ZZ}{\mathbb{Z}}
\newcommand{\pcond}{p_{\mathrm{cond}}}
\newcommand{\pQM}{p_{\mathrm{QM}}}
\newcommand{\xcond}{x_{\mathrm{cond}}}
\newcommand{\CHSH}{\mathrm{CHSH}}
\newcommand{\Tsirelson}{2\sqrt{2}}
\newcommand{\RR}{\mathbb{R}}
\newcommand{\Hil}{\mathcal{H}}
\newcommand{\CQ}{\mathcal{C}_Q}
\newcommand{\CQr}{\mathcal{C}_Q^{\mathrm{red}}}
\newcommand{\bOmega}{\boldsymbol{\Omega}}
\newcommand{\dmu}{d\mu}
\newcommand{\nlvl}[1]{n\!\left(#1\right)}
\newcommand{\IE}{\mathrm{IE}}
\newcommand{\bsF}{{}_2F_1}
\newcommand{\Qgr}{Q_{\mathrm{group}}}
\newcommand{\Mup}{M_{\mathrm{up}}}
\newcommand{\Mdn}{M_{\mathrm{down}}}
\newcommand{\Tw}{\mathrm{Tw}}
\newcommand{\alphas}{\alpha_s}
\newcommand{\alphaem}{\alpha_{\mathrm{em}}}




\begin{theorem}[No $E_7$-type matter sector]
\label{P21:thm:noE7}
If matter sectors are realised as finite subgroups of the vacuum group
$\twoI$, then the binary octahedral group $\twoO$ (McKay $\widehat{E}_7$)
carries no matter sector.
\end{theorem}
\status{draft}
\source{P21:thm:noE7}
\residual{exact}
\begin{proof}
$|\twoO|=48$ and $|\twoI|=120$. Since $48\nmid120$, Lagrange's theorem
forbids $\twoO\subset\twoI$. Equivalently, the octahedral group
$O\cong S_4$ (order $24$) is not a subgroup of the icosahedral group
$I\cong A_5$ (order $60$, with $24\nmid60$), so no lift to the binary
cover exists. Hence no $\widehat{E}_7$ matter sector is available.
\end{proof}
\status{draft}
\source{P21:thm:noE7}

\begin{remark}[The D-type slot is available and colourless]
\label{P21:rem:dtype_slot}
The binary dihedral (dicyclic) groups $\twoD_n$ of order $4n$ are
subgroups of $\twoI$ whenever $4n\mid120$ after the appropriate lift;
in particular $\twoD_3$ (order $12$, $\widehat{D}_5$) and $\twoD_5$
(order $20$, $\widehat{D}_7$) lie in $\twoI$ but \emph{not} in $\twoT$
(since $A_4$ has no subgroup of order $6$, neither dicyclic lifts into
$\twoT$). They are ``sideways branches'' of the vacuum. Crucially,
D-type groups carry their own $\Ztwo$ outer automorphism and not the
colour $\Zthree$ (Paper~XVII, A-type versus E-type McKay chains): they
are colourless, non-baryonic, and reserved for beyond--Standard-Model
physics. The candidate dark sectors are $\twoD_3,\twoD_5$.
\end{remark}
\status{draft}
\source{P21:rem:dtype_slot}

\begin{proposition}[No closing charge for the D-type sector]
\label{P21:prop:noclosure}
Neither natural WZW assignment of the D-type sector admits a closing
integer:
\begin{enumerate}[noitemsep,label=(\alph*)]
\item $(D_{n+2})_1=\SO(2r)_1$ (the McKay-algebra convention, as the quark
uses $(E_6)_1$): for $\twoD_3\to\SO(10)_1$ ($r=5$) and
$\twoD_5\to\SO(14)_1$ ($r=7$), the spinor primary has $h_s=r/8$, forcing
$Q\,r$ odd, while the vacuum $T_1=-r/24$ forces $Q\,r$ even. No integer
$Q$ satisfies both.
\item $\SU(2)_k$ with the D-type modular invariant (the convention the
lepton uses, $k=2n$): all even levels $k$ fail closure, while only odd
$k$ close; the D-type invariants sit at $k=6$ ($\twoD_3$) and $k=10$
($\twoD_5$), both even.
\end{enumerate}
Hence $\Qgr$ does not exist for the D-type sector.
\end{proposition}
\status{draft}
\source{P21:prop:noclosure}
\begin{proof}
Direct computation of the closure condition on the SO$(2r)_1$ primaries
$\{1,v,s,c\}$ with $(h_1,h_v,h_s)=(0,\tfrac12,\tfrac r8)$ and on the
$\SU(2)_k$ primaries $h_\ell=\ell(\ell+2)/4(k+2)$, validated against the
three Standard-Model sectors ($(E_6)_1\to3$, $\SU(2)_1\to6$,
$\SU(2)_3\to10$) before reporting the null D-type results
(\path{src_paper21/dtype_cft_tower_engine.py}).
\end{proof}
\status{draft}
\source{P21:prop:noclosure}

\begin{corollary}[Sterility]
\label{P21:cor:sterile}
T-matrix closure is the charge-quantisation condition; its failure means
the D-type sector carries \emph{no} quantised Standard-Model charge.
Equivalently, electric charge is the $\SU(3)_L$ generator
$Q=T_{3L}+T_{8L}/\sqrt3$ of the $E_6$ trinification (Paper~XV),
quantised in units of $\tfrac13$ by the $\Zthree$ centre of $\SU(3)_L$;
a D-type group carries only its own $\Ztwo$ and no $\Zthree$, so it
admits no $\SU(3)_L$ representation for $Q$ to act on. Being a separate
McKay branch, it has no Standard-Model doublet partner and hence no weak
charge either. The D-type sector is a full Standard-Model singlet:
\emph{sterile}, interacting only gravitationally and through its own
$\Ztwo$.
\end{corollary}
\status{draft}
\source{P21:cor:sterile}

\begin{proposition}[Sterile and Dirac share one origin]
\label{P21:prop:sterile_dirac}
The odd-rank spinor $h_s=r/8$ is simultaneously (i) the primary whose
T-matrix phase obstructs closure (Proposition~\ref{P21:prop:noclosure}a),
making the sector sterile, and (ii) complex ($s^{*}=c$, since both
candidates have $r$ odd: $\SO(10):r=5$, $\SO(14):r=7$), endowing the
sector with a conserved dark-number $U(1)$ and hence a distinct
antiparticle. A complex (Dirac) sector can carry a particle--antiparticle
asymmetry; contrast the neutrino sector $\SU(2)_1$, whose $j=\tfrac12$
primary is self-conjugate ($j^{*}=k-j=\tfrac12$), i.e.\ Majorana
(Paper~XVII), and which therefore cannot. The D-type sector is thus a
\emph{sterile Dirac} sector carrying its own $\Ztwo$ (a dark parity,
stabilising its lightest state) and a complex-spinor $U(1)$ (a dark
number, admitting an asymmetry). Its McKay algebra for $\twoD_3$ is
$\SO(10)$---the archetypal grand-unified group, the home of
baryon-number violation (used here as the McKay-correspondent algebra,
not as a gauged grand-unified theory).
\end{proposition}
\status{draft}
\source{P21:prop:sterile_dirac}

\begin{proposition}[$a^{-3}$ equation of state, conditional]
\label{P21:prop:eos}
Let the sector's gravitating content be a population of localised,
number-conserved coherent configurations of fixed proper size and energy
(solitons, or frozen director deformations held rigid by the condensate's
static shear rigidity, Paper~VII). Then the number density dilutes as
$a^{-3}$ and the energy density as $\rho\propto a^{-3}$, $w=0$: cold
dark matter. If instead the content relaxes to a smooth space-filling
texture, gradients dilute as $a^{-1}$ and $\rho\propto a^{-2}$,
$w=-\tfrac13$: not cold. Which case obtains is set by whether the
configurations stay localised or coarsen; in the framework's inertial
director dynamics, coarsening proceeds at a rate controlled by the
signal speed $\cdir=\sqrt{K/\chi}$, with a coarsening time that diverges
as $\cdir\to0$. A coarse single-quench freeze test
(\path{src_paper7/dm_polarization/defect_network_sim_3d.py}) confirms
the localised ($a^{-3}$) regime for small $\cdir$ and coarsening for
larger $\cdir$, with repeated mergers re-seeding and sustaining the
localised network.
\end{proposition}
\status{draft (conditional: rests on the coldness threshold being met, not derived from first principles)}
\source{P21:prop:eos}

The abundance scenario (asymmetric-dark-matter co-genesis, developed
fully in the companion Paper~XXII) posits one chiral, CP-violating
reconnection event in the shared $\twoI$ vacuum seeding correlated
baryon and dark asymmetries, with branching ratio $r=n_{\mathrm{dark}}/n_b$:
\begin{equation}
  \frac{\Odm}{\Ob}=r\,\frac{m_{\mathrm{dark}}}{m_p},
  \label{P21:eq:adm}
\end{equation}
where the gravitating mass per relic is its rest mass $m_{\mathrm{dark}}$
(the relic is a stable knot; Hopf-charge conservation forbids its
unwinding). Taking the observed $\Odm/\Ob\approx5.4$ with $r\approx1$
(from the comparable subgroup orders $\twoD_3:\twoT=12:24$, not derived)
gives
\begin{equation}
  m_{\mathrm{dark}} \;\approx\; 5\,m_p \;\approx\; 5~\mathrm{GeV},
\end{equation}
the canonical asymmetric-dark-matter mass scale
(\path{src_paper21/adm_abundance.py}). This is a structured,
falsifiable, but explicitly \emph{conditional} scenario: it is not a
derivation of $\Odm$. Four quantities remain unestablished --- the
baryogenesis magnitude $\eta$ (substantially developed in the companion
Paper~XXII, which narrows it to the first-order window's location and
the heavy-band power, both still open), the soliton's existence (its
\emph{stability}, by contrast, is topological and settled), the
branching ratio $r$, and the annihilation efficiency together with the
coldness threshold on $\cdir$ --- enumerated as open problems
P21:OP1--OP4.
\status{draft (conditional scenario, not a derivation of $\Odm$)}
\source{P21:eq:adm}
